\section{Qubit control and readout} \label{sec:control}

\subsection{Control}

Operating the device requires simultaneous synchronized control waveforms for each of the qubits and couplers.
We use 54 coherent microwave control signals for qubit XY rotations, 54 fast flux bias lines for qubit frequency tuning, and 88 fast flux biases for the adjustable couplers.
Dispersive readout requires an additional 9 microwave signals and phase sensitive receivers.
A schematic of the room temperature electronics is shown in Fig.~\ref{fig:sup_control:electronics}, and the cryogenic wiring is shown in Fig.~\ref{fig:sup_control:cryo_wiring}.

Waveform generation is based on a custom-built multi-channel digital to analog converter (DAC) module.
Each DAC module provides 8 DACs with 14-bit resolution and 1 GS/s sample rate.
Each DAC sample clock is synchronized to a global 10 MHz reference oscillator, and their trigger is connected by a daisy chain to synchronize all modules used in the experiment.
This set of DAC modules forms a $>$250-channel, phase-synchronous waveform generator.
We have measured 20 ps of jitter between channels.
The modules are mounted in 14-slot 6U rackmount chassis.
A single chassis, shown in FIG.~\ref{fig:sup_control:rack_photo}, can control approximately 15 qubits including their associated couplers and readout signals.
A total of 4 chassis are used to control the entire Sycamore chip.

The DAC outputs are used directly for fast flux biasing the qubits and couplers required for two-qubit gates.
Microwave control for single-qubit XY rotations and dispersive readout combine two DAC channels and a mixer module to form a microwave arbitrary waveform generator (Microwave AWG) via single-sideband upconversion in an IQ mixer as shown in Figure \ref{fig:sup_control:electronics}\,a.
The microwave AWG provides signals with arbitrary spectral content within $\pm$350 MHz of the local oscillator (LO).
A single LO signal is distributed to all IQ mixers so that all qubits' XY controls are phase coherent.
The mixer modules are mounted in the same chassis as the DAC modules.
Each mixer's I and Q port DC offsets are calibrated for minimum carrier leakage and the I and Q amplitudes and phases are calibrated to maximize image rejection.

Each DAC module contains an FPGA that provides a gigabit ethernet interface, SRAM to store waveform patterns, and sends the waveform data to the DAC module's 8 DACs.
To optimize the use of SRAM, the FPGA implements a simple jump table to allow reusing or repeating waveform segments.
A computer loads the desired waveforms and jump table onto each FPGA using a UDP-based protocol and then requests the first (master) FPGA to start.
The start pulse is passed down the daisy chain causing the remainder (slave) DACs and ADCs to start.

\quickfig{\columnwidth}{figures/sup_control/combined}{\textbf{Control electronics.} \textbf{a,} The custom DAC module provides 8 DAC channels (4 shown). DACs are used individually for flux pulses or in pairs combined with a mixer module to comprise a microwave AWG channel (dashed box). \textbf{b,} A single DAC channel and a microwave source are used to bias and pump the parametric amplifier for readout. \textbf{c,} Readout pulses are generated by a microwave AWG. The reflected signal is amplified, mixed down to IF, and then digitized in a pair of ADCs. The digital samples are analyzed in the FPGA.}{fig:sup_control:electronics}

\quickfig{\columnwidth}{figures/sup_control/cryo_wiring_diagram}{\textbf{Cryogenic wiring.} Control and readout signals are carried to and from the Sycamore chip with a set of cables, filters, attenuators, and amplifiers. }{fig:sup_control:cryo_wiring}

\quickfig{\columnwidth}{figures/sup_control/rack_photo_annotated}{\textbf{Electronics chassis.} Each chassis supports 14 DAC and/or mixer modules. Local oscillators are connected at the top of each mixer module. A set of daisychain cables connects from each ADC module to the next. Control signals exit the chassis through coaxial cables.}{fig:sup_control:rack_photo}

\subsection{Readout}

Qubit state measurement and readout (hereafter ``readout'') are done via the dispersive interaction between the qubit and a far-detuned harmonic resonator \cite{Blais2004CircuitQED, Gambetta2006, bultink2018general}.
A change in the qubit state from $\ket{0}$ to $\ket{1}$ causes a frequency shift of the resonator from $\omega_{\ket{0}}$ to $\omega_{\ket{1}}$.
A readout probe signal applied to the resonator at a frequency in between $\omega_{\ket{0}}$ and $\omega_{\ket{1}}$ reflects with a phase shift $\phi_{\ket{0}}$ or $\phi_{\ket{1}}$ that depends on the resonator frequency and therefore on the qubit state.
By detecting the phase of the reflected probe signal we infer the qubit state.
The readout probe signal is generated with the same microwave AWG as the XY control signals, but with a separate local oscillator, and is received and demodulated by the circuit shown in Figure \ref{fig:sup_control:electronics}\,c.

The readout probe intensity is typically set to populate the readout resonator with only a few photons to avoid readout-induced transitions in the qubit \cite{Sank2016Transitions}.
Detecting this weak signal at room temperature with conventional electronics requires 100 dB of amplification.
To limit the integration time to a small fraction of the qubit coherence time, the amplification chain must operate near the quantum noise limit \cite{clerk2010intro, caves1982Amplifiers}.

Inside the cryostat the signal is amplified by an impedance matched lumped element Josephson parametric amplifier (IMPA) \cite{Mutus2014IMPA} on the mixing chamber stage followed by a Low Noise Factory cryogenic HEMT amplifier at 3 K.
At room temperature the signal is further amplified before it is mixed down with an IQ mixer producing a pair of intermediate frequency (IF) signals $I(t)$ and $Q(t)$.
The IF signals are amplified by a pair of variable gain amplifiers to fine-tune their level, and then digitized by a pair of custom 1 GS/s, 8-bit analog to digital converters (ADC).
The digitized samples $I_n$ and $Q_n$  are processed in an FPGA which combines them into a complex phasor
\begin{equation*}
  z_n = I_n + iQ_n = E_n \exp(i (\omega n dt + \phi))
\end{equation*}
where $dt$ is the sample spacing, $\omega$ is the IF frequency, $\phi$ is the phase that depends on the qubit state, and $E_n$ is the envelope of the reflected readout signal.
The envelope is measured experimentally once and then used by the FPGA in subsequent experiments as the optimal demodulation window $w_n$ to extract the phase of the reflected readout signal \cite{ryan2015tomography, Jeffrey2014Fast}.
The FPGA multiplies $z_n$ by $w_n \exp(-i \omega n dt)$, and then sums over time to produce a final complex value $\exp(i \phi)$
\begin{equation*}
  \sum_{n=0}^{N-1} z_n w_n \exp(-i \omega n dt) \propto \exp(i \phi)
\end{equation*}
In the absence of noise, the final complex value would always be one of two possible values corresponding to the qubit states $\ket{0}$ and $\ket{1}$.
However, the noise leads to Gaussian distributions centered at those two points.
The size of the clouds is determined mostly by the noise of the IMPA and cryogenic HEMT amplifier, while the separation between the clouds' centers is determined by the resonator probe power and duration.
The signal to noise ratio of the measurement is determined by the clouds' separation and width \cite{Jeffrey2014Fast, Sank2015Noise}.

The 54 qubits are divided into nine frequency multiplexed readout groups of six qubits each.
Within a group, each qubit is coupled to its own readout resonator, but all six resonators are coupled to a shared bandpass Purcell filter \cite{Reed2010Purcell, Sete2015Purcell, Jeffrey2014Fast}.
All qubits in a group can be read-out simultaneously by frequency-domain multiplexing \cite{chen2012mux, barends2014nine} in which the total probe signal is a superposition of probe signals at each of the readout resonators’ frequencies.
The phase shifts of these superposed signals are independently recovered in the FPGA by demodulating the complex IQ phasor with each intermediate frequency.
In other words, we know what frequencies are in the superposed readout signal and we compute the Fourier coefficients at those frequencies to find the phase of each reflected frequency component.
